%% ---------------------------------------------------------------------
\section{Un ejemplo para verlas actuar}
\label{sec:pn-ejemplo}
%% ---------------------------------------------------------------------

\begin{ejemplo}[label=ej:pn-mra]{Masa, resorte y amortiguador}
El sistema $M\ddot x+b\dot x+kx=F$, con $M=1$\,kg, $b=10$\,N\,s/m y $k=20$\,N/m, tiene
\[
  \Gp(s)=\frac{X(s)}{F(s)}=\frac{1}{s^2+10s+20}.
\]
Se quiere, con retroalimentación unitaria y un PID en forma paralelo, un tiempo de subida corto, el menor sobrepico posible y error cero ante escalón.
\solucion \textbf{Lazo abierto.} Los polos son $-2.76$ y $-7.24$ ($\zeta=1.12$): sin sobrepico, con $\tr=0.88$\,s, $\ts=1.60$\,s y valor final $1/20$. Respecto a una referencia unitaria, el error es del $95\%$.

\textbf{P.} El lazo es $s^2+10s+(20+\Kc)$: $\wn=\sqrt{20+\Kc}$, $\zeta=5/\sqrt{20+\Kc}$ y $\ess=20/(20+\Kc)$. Subir $\Kc$ reduce el tiempo de subida y el error, pero baja $\zeta$ y sube el sobrepico. Con $\Kc=300$: $\ess=0.0625$, $\zeta=0.28$ y $\OS=40\%$ (\cref{fig:pn-mra}, arriba a la izquierda).

\textbf{PD.} Con $\Kc=300$, la derivada agrega amortiguamiento: el lazo es $s^2+(10+K_d)s+320$. Con $K_d=20$, $\zeta=0.84$ y el sobrepico cae al $5.8\%$; el error no cambia ($K_ds$ se anula en $s=0$). El PD introduce además un cero en $-\Kc/K_d=-15$; por eso el sobrepico real es mayor que el que predice el segundo orden sin cero ($0.8\%$).

\textbf{PI.} Con $\Kc=300$, la integral anula el error, pero el lazo es de tercer orden, $s^3+10s^2+320s+K_i$, estable solo si $K_i<3200$ (Routh). Con $K_i=200$: $\ess=0$ y $\OS=36\%$.

\textbf{PID.} Con $\Kc=350$, $K_i=300$ y $K_d=50$, la respuesta no tiene sobrepico, sube en $0.05$\,s, se establece en $0.83$\,s y no tiene error. Los ceros del controlador ($-1$ y $-6$) casi cancelan los polos más lentos del lazo cerrado.
\end{ejemplo}

\begin{figure}[ht]
\centering
\includegraphics[width=\linewidth]{cap07/f_mra}
\caption{Ejemplo~\ref{ej:pn-mra}: efecto de cada acción sobre la respuesta al escalón.}
\label{fig:pn-mra}
\end{figure}

\codigorepo{cap07_pid_nociones/C7_Ej2_Acciones.m}{pid}{ejemplo~\ref{ej:pn-mra}: P, PD, PI y PID sobre el sistema masa-resorte-amortiguador}

\begin{table}[ht]
\centering
\caption{Efecto típico de aumentar cada parámetro de un PID sobre una planta estable~\citep{ACL05}. Son tendencias, no leyes: el problema~\ref{prob:pn-mra} pide comprobarlas.}
\label{tab:pn-efectos}
\small
\begin{tabular}{@{}lccccc@{}}
\toprule
Aumentar & $\tr$ & $\OS$ & $\ts$ & $\ess$ & Estabilidad\\
\midrule
$\Kc$ & baja & sube & sube un poco & baja & empeora\\
$K_i=\Kc/\Ti$ & baja un poco & sube & sube & se anula & empeora\\
$K_d=\Kc\Td$ & baja un poco & baja & baja & casi no cambia & mejora (si no hay ruido)\\
\bottomrule
\end{tabular}
\end{table}

\begin{receta}[Ajuste manual de un PID]
\begin{enumerate}
  \item Mire la respuesta en lazo abierto y decida qué hay que mejorar.
  \item Agregue acción proporcional para acelerar la respuesta.
  \item Agregue acción derivativa si el sobrepico es excesivo.
  \item Agregue acción integral para eliminar el error en estado estacionario.
  \item Ajuste las tres hasta lograr la respuesta deseada, y verifique la estabilidad.
  \item No use las tres si no hace falta: la derivada trae problemas (\cref{sec:pn-modificaciones}). Mantenga el controlador tan simple como se pueda.
\end{enumerate}
\end{receta}

En resumen: el \textbf{P} acelera pero deja error (salvo en procesos integradores); el \textbf{PI} elimina el error pero oscila más y pierde estabilidad al subir $\Kc$; el \textbf{PD} anticipa y amortigua, pero deja el error del P; el \textbf{PID} junta las ventajas, a costa de amplificar el ruido. En el capítulo~10 (\cref{sec:rf-pid}) veremos lo mismo en frecuencia: un PI es un compensador de atraso, un PD uno de adelanto y un PID uno de atraso-adelanto.
